\documentclass[11pt,a4paper]{article}
\usepackage[utf8]{inputenc}
\usepackage[T1]{fontenc}
\usepackage{lmodern}
\usepackage{amsmath,amssymb}
\usepackage[colorlinks=true,linkcolor=blue,citecolor=blue,urlcolor=blue]{hyperref}
\usepackage{geometry}
\usepackage[protrusion=true,expansion=false]{microtype}
\geometry{margin=2.8cm}

\title{\textbf{Related Work}\\[0.8ex]
\large Principia: A Real-Time WebGPU Explorer\\
of the Three-Body Initial-Condition Manifold}
\author{Malachy Doherty}
\date{March 2026}

\begin{document}
\maketitle
\tableofcontents
\vspace{1em}

% =========================================================
\section{Introduction}
% =========================================================

Principia is a browser-based instrument for exploring the initial-condition (IC) manifold of
the planar gravitational three-body problem in real time.  It renders a zoomable
two-dimensional colour map whose every pixel represents one symmetry-reduced IC, and whose
colour encodes a per-trajectory diagnostic computed on the GPU.  The Burrau--Pythagorean
research programme uses Principia as a computational platform for a systematic survey of the
infinite family of configurations indexed by primitive Pythagorean triples, together with a
continuous extension into mass and momentum space around each member of that family.

This document surveys the research traditions on which Principia draws---IC-manifold mapping,
shape-space geometry, symbolic dynamics of triple encounters, chaos diagnostics, statistical
theory of outcomes, symplectic integration, and GPU-accelerated computation---and identifies
the gaps in the existing literature that motivate the present work.

This document also identifies, in the gaps section and where noted explicitly, several
constructs proposed as part of the present work rather than as established literature: a
free-group formulation of encounter sequences (\S\ref{sec:symbolic}), a ``projective
microscope'' interaction mode for rotating the 2D rendering plane within the 8D manifold
(\S\ref{sec:number_theory}), and the use of the Pythagorean-triple family as a structured
indexing scheme for three-body IC space (\S\ref{sec:number_theory}).  These are clearly
labelled as proposed contributions wherever they appear.

A note on the object of study is useful at the outset.  After fixing the conserved quantities
of the planar three-body problem (total energy, total angular momentum, total mass) and quotienting
by Euclidean symmetry, the space of distinct initial conditions is 8-dimensional: two configuration
coordinates, four momentum coordinates, and two mass-ratio coordinates.  All prior pixel-scale IC
surveys are restricted to a 2-dimensional slice of this space (equal masses, zero velocities).
Principia is motivated by the need to navigate \emph{families} of such 2D slices inside the full
8-dimensional manifold, making the structure of that manifold---not merely one or two cross-sections
of it---the primary object of study.

% =========================================================
\section{The Three-Body Problem: Fundamentals and History}
% =========================================================

The three-body problem---determining the motion of three mutually gravitating point masses---was
formulated in its modern form by Newton in the \emph{Principia Mathematica} of
1687~\cite{newton1687}.  Despite the simplicity of the governing equations, two exceptional
families of closed-form solutions were found by Euler (collinear configurations, 1763) and
Lagrange (equilateral triangle configurations, 1772), but no general analytic solution is known.

The non-existence of a third independent first integral beyond energy and angular momentum was
established by Bruns~\cite{bruns1887} and, more penetratingly, by Poincar\'{e}, whose
\emph{M\'{e}thodes Nouvelles de la M\'{e}canique C\'{e}leste}~\cite{poincare1892} showed that
the problem is generically non-integrable: small perturbations of stable periodic solutions lead
to homoclinic tangles and chaotic behaviour, implying that the IC space must contain fractal
basin boundaries.

A formal power-series solution converging for all time was given by Sundman~\cite{sundman1912},
but the series converges too slowly for practical computation.  The problem must therefore be
treated numerically.  Comprehensive introductions from the celestial-mechanics perspective are
given by Szebehely~\cite{szebehely1967book}, Valtonen et al.~\cite{valtonen2016}, and Heggie
\& Hut~\cite{heggie2003}.

Generic bound three-body systems almost always dissolve into a stable binary and an escaping
single body~\cite{heggie2003}; understanding how the outcome depends on initial conditions, and
charting the fractal structure of that dependence across the full IC manifold, is the central
problem that the work described here addresses.

% =========================================================
\section{The Pythagorean Three-Body Problem}
% =========================================================

The configuration now called the \emph{Pythagorean three-body problem} was introduced by
Burrau in 1913~\cite{burrau1913}: three bodies with masses $m_1=3$, $m_2=4$, $m_3=5$ placed
at rest at the vertices of a 3--4--5 right triangle, with each mass equal to the length of the
opposite side.  Szebehely \& Peters~\cite{szebehely1967} gave the first complete numerical
solution using regularisation: after a protracted series of close triple approaches, the
lightest body escapes in a hyperbolic orbit, leaving a bound binary of the two heavier bodies.
This trajectory has since been a canonical benchmark for three-body codes and a standard
demonstration of sensitivity to initial conditions in celestial mechanics.

Aarseth, Anosova, Orlov \& Szebehely~\cite{aarseth1994} investigated the chaoticity structure
near the standard Pythagorean IC, sampling points on concentric circles of radii $5 \times
10^{-8}$ to $5 \times 10^{-3}$ around each body's position and applying time-reversal accuracy
tests.  They identified regular regions adjacent to regions where the escape angle and the
identity of the escaping body vary abruptly with IC---a first map of the IC-space fractal near
the Pythagorean point.

The (3,4,5) configuration is the simplest member of an infinite family indexed by primitive
Pythagorean triples: every primitive triple $(a,b,c)$ with $a^2+b^2=c^2$ yields a distinct
right-triangle geometry, and the Burrau mass convention (mass equal to the opposite side,
normalised) yields a distinct three-body configuration with rationally related mass ratios.
The number-theoretic structure of this family and its implications for the research programme
are discussed in Section~\ref{sec:number_theory}.

% =========================================================
\section{Mapping the Initial-Condition Manifold}
% =========================================================

\subsection{The Agekian--Anosova map}

The earliest systematic attempt to classify the IC space of the general three-body problem is
due to Agekian \& Anosova~\cite{agekian1967}.  They observed that any equal-mass, zero-velocity
planar three-body configuration is determined (up to overall scale and rotation) by the shape of
the triangle formed by the three bodies.  By fixing the two most-separated bodies on a
horizontal axis and allowing the third to vary, all non-equivalent ICs map to a bounded
``curved triangle'' region in the plane---now called the Agekian--Anosova map.  They divided
this region into four qualitatively distinct sub-regions (Lagrangian L, Middle M, Aligned A,
Hierarchical H) and sampled each, establishing a statistical picture of the final-state
distribution.  This work introduced the paradigm of treating the IC space as a two-dimensional
map to be surveyed by colour-coding the outcome, and it underpins all subsequent
IC-manifold mapping work.

A key limitation of the Agekian--Anosova map is that it fixes two degrees of freedom: equal
masses and zero initial velocities.  These are special conditions; the equal-mass, zero-velocity
surface is a measure-zero subset---a codimension-6 submanifold---of the full 8-dimensional
symmetry-reduced IC space for the planar three-body problem.  The dynamics in the remaining six
dimensions (independent mass variation and non-zero momenta) are not represented.

\subsection{The Lehto et al.\ pixel map and the Maasai shield}

Lehto, Kotiranta, Valtonen, Heinäm\"{a}ki, Mikkola \& Chernin~\cite{lehto2008} carried out a
pixel-by-pixel computation over the Agekian--Anosova domain using approximately 200,000
simulations, colour-coding each IC by the lifetime of the triple system before dissolution.
The resulting image---symmetric about both axes due to the equal-mass, zero-angular-momentum
constraint---resembles a Maasai warrior's shield: broad bands of short-lived systems flanked by
intricate swirling structures of longer-lived systems.  Zooming in reveals the same band
structure repeating at every scale, confirming the fractal character of the basin boundaries.
In related work, Heinäm\"{a}ki, Lehto, Valtonen \& Chernin~\cite{heinamaki1998} analysed the
three-body flow using intermittency and found a strange attractor with Hausdorff dimension
$D \approx 2.1$ in the dynamical flow (not in the 2D IC map, but in the higher-dimensional
phase-space trajectory).

This study is the closest existing precedent to Principia's rendering task.  It is, however,
subject to the same restrictions as the Agekian--Anosova map: equal masses, zero velocities,
and a single diagnostic (lifetime).  The computation was performed in batch over an extended
period, with no provision for interactive navigation or zoom.  The full momentum space
orthogonal to the zero-velocity surface remains unexplored by this approach.  For Principia,
the Lehto et al.\ study both validates the rendering paradigm and defines the narrow 2D subspace
within which all prior high-resolution work has been done.

\subsection{Trani et al.\ and isles of regularity}

Trani et al.~\cite{trani2024} recently performed a large-scale IC scan using the \emph{Tsunami}
regularised N-body code, running $\mathcal{O}(10^6)$ scattering encounter simulations over a
2D parameterisation that varies the phase of a pre-existing binary and the approach angle of a
third body.  They found that regular (non-chaotic) trajectories occupy between 28\% and 84\% of
the scanned IC space, forming connected ``isles of regularity'' in the chaotic sea.  This
confirms the picture of a fractal IC-space boundary interspersed with structured regular
regions and demonstrates that modern hardware can sustain IC surveys at million-trajectory scale.

The scattering parameterisation used by Trani et al.\ is specifically adapted to hierarchical
encounter configurations (one body arriving from far away to interact with a pre-existing
binary) and does not cover the non-hierarchical regime in which all three bodies start with
comparable separations.  It is also static: results are computed offline and cannot be explored
interactively.

\subsection{Fractal basin boundaries}

The theoretical framework for fractal basin boundaries in dynamical systems was established by
McDonald, Grebogi, Ott \& Yorke~\cite{mcdonald1985}, who showed that the boundary between
basins of attraction of competing attractors is generically fractal when the underlying dynamics
are chaotic.  Aguirre, Vallejo \& Sanju\'{a}n~\cite{aguirre2001} further showed that in open
Hamiltonian systems with three or more escape channels---as the three-body problem has---the
basin boundaries often satisfy the stronger \emph{Wada} property: every boundary point is a
limit point of three or more distinct basins simultaneously, making the topology of the
boundary dramatically more complex than a simple fractal curve.  For the three-body problem
the relevant escape channels are which body escapes and in which direction; the basin
boundaries inherit this Wada topology from the underlying homoclinic structure.

The KAM theorem (Kolmogorov~\cite{kolmogorov1954}, Arnold~\cite{arnold1963},
Moser~\cite{moser1962}) guarantees the persistence of regular, non-chaotic islands in
nearly-integrable systems.  The three-body problem is not nearly integrable, but KAM-like
regular islands are nonetheless observed in numerical surveys~\cite{trani2024}, consistent with
the general picture of a fractal boundary between chaotic and regular regions.

% =========================================================
\section{Shape Space and Symmetry Reduction}
% =========================================================

\subsection{Jacobi coordinates}

The standard reduction from Cartesian coordinates to a system adapted to the three-body
structure uses Jacobi coordinates~\cite{jacobi1842}: the relative vector $\boldsymbol{\rho}$
between an ``inner pair'' of bodies, and the vector $\boldsymbol{\lambda}$ from the inner-pair
barycentre to the third body.  These four two-dimensional vectors, together with the decoupled
centre-of-mass motion, fully describe the planar three-body system.  Jacobi coordinates are
the standard choice in both analytical and numerical three-body work, from classical perturbation
theory~\cite{poincare1892} to modern N-body codes~\cite{rein2012}.

\subsection{Shape space and the shape sphere}

Fixing the overall scale and orientation of the triangle formed by the three bodies leaves a
two-dimensional space of triangle \emph{shapes}---the \emph{shape sphere} $S^2$.
Montgomery~\cite{montgomery2015} gives a particularly elegant exposition: the shape sphere is
the unit sphere in a natural Euclidean space, with the Lagrange equilateral configurations at
the north and south poles, collinear configurations on the equator, and the shape-sphere metric
related to the mass-weighted kinematic metric on configuration space.  The three pairwise
binary-collision loci correspond to three great circles on the sphere, dividing it into eight
regions.  Choreographic solutions such as the figure-eight~\cite{chenciner2000} trace closed
curves on the shape sphere, and their structure is most naturally understood in this
representation.

The shape sphere provides a natural coordinate system for the zero-velocity IC space: every
point on $S^2$ corresponds to a triangle shape and, together with a scale and a mass
assignment, determines a full zero-velocity IC.  The Agekian--Anosova map is, in this language,
a chart on the shape sphere restricted to equal masses.  For Principia, the shape sphere
supplies the canonical two-dimensional configuration chart used in the rendering pipeline, and
its geometric structure (collision loci, symmetry axes, choreography curves) guides the design
of the chart parameterisations.

\subsection{Symmetry reduction}

Moeckel~\cite{moeckel2023} gives a modern treatment of the natural dynamical reduction of the
three-body problem, systematically factoring out the Euclidean symmetry group (translations,
rotations) and the scaling freedom of Newtonian gravity to obtain a reduced system.  For the
planar problem, after additionally fixing the conserved angular momentum and total mass, the
symmetry-reduced IC space has 8 degrees of freedom: 2 configuration coordinates (e.g.\ the
shape-sphere angles), 4 momentum coordinates, and 2 mass-ratio coordinates.  This 8-dimensional
space is the natural domain for any complete survey of three-body IC-manifold structure.  All
prior IC maps---including the Agekian--Anosova map and the Lehto et al.\ pixel map---are
restricted to the 2-dimensional zero-velocity, equal-mass subspace, which is a
codimension-6 surface within this 8-dimensional manifold.

% =========================================================
\section{Symbolic Dynamics of Triple Encounters}
\label{sec:symbolic}
% =========================================================

\subsection{Classical symbolic dynamics}

Symbolic dynamics assigns a finite alphabet to the regions of phase space and records the
sequence of symbols as a trajectory crosses region boundaries, reducing the continuous dynamics
to a discrete combinatorial object.  The approach was introduced to celestial mechanics by
Alekseev~\cite{alekseev1969}, who applied it to the Sitnikov problem (a massless body
oscillating through the plane of two equal-mass bodies in circular orbit) and proved the
existence of quasirandom trajectories: for any symbol sequence over a two-symbol alphabet,
there exists a trajectory realising that sequence.  This establishes the full topological
complexity of the symbolic dynamics for a gravitational system.

\subsection{Regional coding of the three-body encounter}

In the context of the general three-body problem, Agekian \& Anosova's HAML decomposition
provides a coarse four-symbol alphabet for coding trajectories.  Heinäm\"{a}ki, Lehto, Valtonen
\& Chernin~\cite{chernin2000} used sequences of region visits to define and compute the
Kolmogorov--Sinai entropy of the IC space (see also Section~\ref{sec:ksentropy}).  This is
an early example of using a symbolic sequence to produce a per-IC scalar field over the
Agekian--Anosova map.

\subsection{Tanikawa--Mikkola symbolic dynamics}

Tanikawa \& Mikkola~\cite{tanikawa2015,tanikawa2000} systematically associated symbol sequences
with orbits in the free-fall (zero-velocity) three-body domain, using the equator of the shape
sphere as a natural symbol boundary.  Their alphabet records which binary forms at each close
approach, yielding sequences that characterise periodic and collision orbits.  This work
establishes symbol sequences over the shape sphere as a productive tool for cataloguing
three-body orbits, and it provides a concrete foundation for more refined symbolic schemes.

\paragraph{Proposed contribution.}
In contrast, the present work proposes equipping the encounter alphabet with orientation
(prograde or retrograde) and the algebraic structure of a free group, so that encounter
sequences satisfy composition laws under trajectory concatenation.  Such a free-group
formulation is more sensitive to the signed topology of the encounter sequence than the
regional codings of Agekian--Anosova or the unsigned sequences of Tanikawa--Mikkola; it is
not established in the prior literature.

\subsection{Heggie's encounter classification}

The binary--single scattering problem has been systematically studied by Heggie~\cite{heggie1975}
and Heggie \& Hut~\cite{heggie2003}.  They classify outcomes as fly-by (the single body passes
without disrupting the binary), exchange (the single replaces one binary member), and
ionisation.  These coarse outcome classes correspond to distinct regions in the IC space of the
scattering problem, and the boundaries between them are the three-body analogues of the fractal
basin boundaries visible in the Agekian--Anosova and Lehto maps.

% =========================================================
\section{Chaos Diagnostics}
\label{sec:chaos}
% =========================================================

\subsection{Lyapunov exponents and FTLE}

The primary quantitative diagnostic of chaos in Hamiltonian systems is the Lyapunov
characteristic exponent, measuring the exponential rate of divergence of nearby trajectories.
Benettin, Galgani, Giorgilli \& Strelcyn~\cite{benettin1980} established the theoretical
foundations and gave a practical algorithm for computing the full Lyapunov spectrum via periodic
Gram--Schmidt reorthonormalisation of a tangent-vector basis evolved by the linearised flow.

For comparing chaos across different ICs over a finite integration horizon, the
\emph{finite-time Lyapunov exponent} (FTLE) is more appropriate than the asymptotic exponent.
Haller~\cite{haller2001} established the FTLE field---computed over a spatial grid of initial
conditions---as the standard tool for identifying Lagrangian coherent structures in fluid
dynamics, where the ridges of the FTLE field delineate the boundaries of qualitatively different
flow regimes.  Applied to the three-body IC space, an FTLE field rendered as a colour map over
the IC plane delineates the fractal basin boundaries more sharply than a simple escape-outcome
map, because it is sensitive to trajectory divergence rather than just final state.  We have not
found prior work that computes the FTLE as a 2D field over a dense pixel grid for the three-body
IC space; the present work applies this tool in that setting.

\subsection{Frequency diffusion and spectral chaos indicators}
\label{sec:freqdiff}

An independent chaos indicator is the \emph{frequency diffusion coefficient}, which measures
the drift of fundamental orbital frequencies across sub-intervals of a trajectory.  Laskar
introduced frequency analysis as a tool for identifying regular and chaotic orbits in the solar
system~\cite{laskar1990,laskar1993}: on a KAM torus the fundamental frequencies are conserved;
in the chaotic region they diffuse, and the diffusion rate provides a chaos indicator more
sensitive than the short-time Lyapunov exponent.  The relationship between frequency diffusion
and Arnold diffusion in multi-dimensional Hamiltonian systems is discussed by
Nekhoroshev~\cite{nekhoroshev1977} and Laskar and collaborators.

Frequency diffusion has been applied extensively in the solar system context but not, to our
knowledge, rendered as a 2D field over a dense IC grid for the general three-body problem.

\subsection{Kolmogorov--Sinai entropy in the three-body problem}
\label{sec:ksentropy}

The Kolmogorov--Sinai (KS) entropy---the rate of information creation along a trajectory,
equal to the sum of positive Lyapunov exponents---provides a global chaos measure per IC.
Heinäm\"{a}ki et al.~\cite{chernin2000} applied KS-entropy analysis to the
Agekian--Anosova IC space, producing entropy maps that correlate strongly with the lifetime
maps of Lehto et al.~\cite{lehto2008}: short-lived systems tend to low entropy, long-lived
systems to high entropy.  Chernin, Valtonen \& Mikkola~\cite{chernin1994} carried out a
related KS-entropy analysis for an analogous restricted four-body system.  These entropy maps
are early examples of multi-diagnostic rendering of IC-space structure---computing something
other than mere lifetime---but are restricted to the equal-mass, zero-velocity subspace (or
to an analogous restricted configuration).

% =========================================================
\section{Statistical Theory of Three-Body Outcomes}
% =========================================================

Because individual three-body trajectories are unpredictable in the chaotic regime, a
complementary approach characterises the \emph{distribution} of outcomes over ensembles of ICs.
Heggie~\cite{heggie1975} derived the first statistical theory of binary-single scattering, using
a microcanonical ensemble argument to predict post-encounter binary orbital element distributions.
This theory was extended by Hut \& Bahcall~\cite{hutbahcall1983} and applied broadly in the
context of dense stellar systems.

Stone \& Leigh~\cite{stone2019} gave a modern, closed-form statistical solution to the
non-hierarchical three-body problem under the ergodicity assumption: the distribution of
outcomes agrees well with numerical ensembles for ``resonant'' encounters (those involving
multiple close triple approaches, approximately 50\% of all encounters~\cite{stone2019}) but fails for the
non-resonant (``prompt'') fraction, which escapes without full ergodic mixing.

Breen et al.~\cite{breen2020} trained a deep neural network to predict three-body trajectories
at a fraction of the cost of direct integration, demonstrating that the IC space contains
learnable structure; however, accuracy degrades near fractal basin boundaries, exactly the
regime where Principia's diagnostic rendering is most informative.

The tension between the Stone--Leigh statistical theory and the structured, non-ergodic
features visible in IC maps such as the Lehto et al.\ Maasai shield is an open question: the statistical theory
predicts smooth distributions, while the IC maps reveal sharp fractal boundaries that imply
strongly non-ergodic structure at fine scales.  Dense IC maps with multiple simultaneous
diagnostics are a natural tool for investigating where and how the ergodic approximation breaks
down as a function of IC.  For Principia, this motivates computing both chaos-sensitive
diagnostics (FTLE, frequency diffusion) and outcome-based diagnostics (escape identity, orbit
count) simultaneously, so that the two can be compared directly within the same IC neighbourhood.

% =========================================================
\section{Symplectic Numerical Integration}
% =========================================================

Symplectic integrators, which exactly preserve the symplectic structure of Hamiltonian systems,
maintain bounded energy error over exponentially long times, making them the standard choice for
long-duration celestial mechanics simulations~\cite{leimkuhler2004}.  The standard textbook
treatment is Leimkuhler \& Reich~\cite{leimkuhler2004}; applications to N-body celestial
mechanics are discussed by Heggie \& Hut~\cite{heggie2003}.

Yoshida~\cite{yoshida1990} showed how to construct symplectic integrators of arbitrarily high
order by composing leapfrog steps with appropriately chosen step sizes, giving 4th and 6th
order methods whose energy-conservation properties make them well suited to long IC-manifold
scans in which each trajectory must be integrated to a common horizon.

For close encounters between bodies---triple approaches and near-collisions---symplectic
integrators at fixed step size lose accuracy.  Aarseth~\cite{aarseth2003} developed
regularisation methods (Kustaanheimo--Stiefel and its extensions) that analytically resolve
close pairs.  Modern codes such as REBOUND~\cite{rein2012} and Tsunami~\cite{trani2024} use
variants of these approaches for close-encounter handling; their treatment of this regime is
more accurate than a fixed-step integrator, at the cost of variable computational load per
trajectory.  Li \& Liao~\cite{li2017} additionally demonstrated that reliable long-time
trajectories in the chaotic three-body problem require extended-precision arithmetic (their
``clean numerical simulation''), identifying numerical noise as a fundamental constraint on
trajectory fidelity in the deeply chaotic regime.

Principia adopts a fixed-step Yoshida leapfrog on the GPU for parallelism and uniformity.
This means close-encounter trajectories are not regularised; near-collision samples are
flagged and the trajectory is terminated rather than continued.  This bounds Principia's
reliable IC-map coverage to initial conditions that do not involve deep triple encounters,
and means that certain diagnostics (orbit count, encounter word) may be unavailable or
unreliable near collision loci.  These limitations are a known trade-off of the GPU
parallelism architecture.

% =========================================================
\section{GPU-Accelerated Parallel Computation}
% =========================================================

Computing an IC map at $N^2$ pixels, each requiring an independent trajectory integration, is
embarrassingly parallel: no data is shared between pixels.  This maps naturally to the massively
parallel architecture of modern graphics processing units.

GPU acceleration for gravitational N-body problems was pioneered by Portegies Zwart and
collaborators~\cite{portegies2007}, achieving order-of-magnitude speedups over CPU codes for
direct-sum N-body.  The REBOUND code~\cite{rein2012} provides a widely used CPU-based N-body
framework with symplectic integrators; it does not, however, support GPU execution.

The WebGPU API~\cite{webgpu2024} is a new W3C standard that exposes GPU compute hardware
through a browser-based sandbox, allowing compute-shader programmes to run without native
binary installation.  To our knowledge, the application of WebGPU to gravitational N-body
simulation or interactive IC-manifold visualisation has not been previously reported.

% =========================================================
\section{The Burrau--Pythagorean Family and Number Theory}
\label{sec:number_theory}
% =========================================================

The number-theoretic object at the heart of the Burrau--Pythagorean programme is the set of
primitive Pythagorean triples, parametrised by Euclid's formula~\cite{hardy1979}:
$a = m^2-n^2$, $b = 2mn$, $c = m^2+n^2$ for coprime $m>n>0$ with $m-n$ odd.  This gives
approximately $\mathcal{O}(N/\log N)$ primitive triples with hypotenuse $c \leq N$.  Each
triple defines a distinct Burrau initial condition; the full family sweeps out a 1D curve in
the 8-dimensional symmetry-reduced IC manifold as the Euclid parameters $(m,n)$ vary.

The literature on periodic orbits of the three-body problem provides an important reference
point.  \v{S}uvakov \& Dmitra\v{s}inovi\'{c}~\cite{suvakov2013} reported 13 new families of
planar equal-mass periodic orbits using a topological classification based on shape-sphere
choreography curves; this tripled the number of known periodic families overnight.  Li \&
Liao~\cite{li2017} subsequently used high-precision ``clean numerical simulation'' to find
more than 600 additional families.  These periodic-orbit catalogues are relevant because
Principia's rendering pipeline may reveal periodic-orbit loci as distinguished islands of
regularity within the IC maps, and because the shape-sphere paths of known choreographies
provide reference curves against which computed trajectories can be compared.

We have not found prior dynamical-systems work that uses the Pythagorean-triple family as an
indexing scheme for three-body configurations; the number-theoretic structure motivating the
programme is new to this context.

% =========================================================
\section{Summary of Gaps}
% =========================================================

The literature reviewed above establishes a clear lineage from Agekian \& Anosova (1967) through
Lehto et al.\ (2008) and Trani et al.\ (2024): the IC manifold of the three-body problem can be
rendered as a colour-coded 2D image, revealing fractal basin boundaries, regular islands, and
entropy structure.  The theoretical foundations in shape space, symbolic dynamics, Lyapunov
theory, and statistical mechanics are well developed.  The following gaps are identified.

\paragraph{Coverage of the IC manifold.}
All prior pixel-scale IC surveys are restricted to a 2-dimensional subspace of the
8-dimensional symmetry-reduced IC manifold.  Specifically, the Agekian--Anosova tradition fixes
equal masses and zero velocities; Trani et al.\ fix a hierarchical encounter geometry.  The four
momentum dimensions and the two mass dimensions are entirely absent from existing maps.

\paragraph{Interactivity and zoom.}
All prior IC maps were computed offline in batch.  We have not identified a tool that permits
interactive navigation---zoom, pan, chart switching---within a live-rendered IC map.

\paragraph{Simultaneous multi-diagnostic output.}
Prior maps render a single diagnostic per computation (usually lifetime or escape identity).
We have not found prior work that computes FTLE, frequency diffusion, orbit count, and a
symbolic encounter summary simultaneously over a dense IC grid.

\paragraph{Symbolic encounter invariants.}
Existing symbolic dynamics work uses coarse regional codings (Agekian--Anosova) or unsigned
shape-sphere sequences (Tanikawa--Mikkola).  A group-theoretic formulation that records the
signed (prograde/retrograde) topology of the encounter sequence and satisfies composition laws
has not, to our knowledge, been reported.

\paragraph{Variable-mass IC maps.}
We have not found IC-map work that includes mass ratio as a free parameter or chart axis.
The effect of mass variation on basin topology, fractal dimension, and regular-island structure
is not systematically documented in the literature we have reviewed.

\paragraph{Structured family surveys.}
We are not aware of a systematic survey of how three-body dynamics varies across the
Burrau--Pythagorean family, or any other number-theoretically indexed family of configurations.

\paragraph{Scope of the review.}
This review is necessarily incomplete; the gaps identified are those apparent from the
literature accessible to us at the time of writing.  It is possible that unpublished code,
conference proceedings, or adjacent-field work addresses one or more of these points, and any
strong gap claim should be treated as a research hypothesis rather than an established fact.
Additionally, even with Principia, interactive 2D slice navigation does not equal exhaustive
coverage of the 8-dimensional IC manifold: the tool makes selected 2D cross-sections navigable,
and the structure of the full manifold can only be inferred indirectly from families of such
cross-sections.

% =========================================================
\begin{thebibliography}{99}
% =========================================================

\bibitem{aguirre2001}
Aguirre, J., Vallejo, J.\,C., \& Sanju\'{a}n, M.\,A.\,F. (2001).
Wada basins and chaotic invariant sets in the H\'{e}non--Heiles system.
\textit{Physical Review E}, \textbf{64}, 066208.

\bibitem{aarseth1994}
Aarseth, S.\,J., Anosova, J.\,P., Orlov, V.\,V., \& Szebehely, V.\,G. (1994).
Global chaoticity in the Pythagorean three-body problem.
\textit{Celestial Mechanics and Dynamical Astronomy}, \textbf{58}, 1--16.

\bibitem{aarseth2003}
Aarseth, S.\,J. (2003).
\textit{Gravitational N-Body Simulations}.
Cambridge University Press.

\bibitem{agekian1967}
Agekian, T.\,A., \& Anosova, Z.\,P. (1967).
A study of the dynamics of triple systems by means of statistical sampling.
\textit{Astrophysics}, \textbf{3}, 1--17.

\bibitem{alekseev1969}
Alekseev, V.\,M. (1969).
Quasirandom dynamical systems, I--III.
\textit{Mathematics of the USSR-Sbornik}, \textbf{5--7} (series of papers, 1968--1969).

\bibitem{arnold1963}
Arnold, V.\,I. (1963).
Proof of a theorem of A.\,N.\,Kolmogorov on the preservation of conditionally periodic motions.
\textit{Russian Mathematical Surveys}, \textbf{18}(5), 9--36.

\bibitem{benettin1980}
Benettin, G., Galgani, L., Giorgilli, A., \& Strelcyn, J.-M. (1980).
Lyapunov characteristic exponents for smooth dynamical systems and for Hamiltonian systems:
a method for computing all of them.
\textit{Meccanica}, \textbf{15}, 9--30.

\bibitem{breen2020}
Breen, P.\,G., Foley, C.\,N., Boekholt, T., \& Portegies Zwart, S. (2020).
Newton versus the machine: solving the chaotic three-body problem using deep neural networks.
\textit{Monthly Notices of the Royal Astronomical Society}, \textbf{494}, 2465--2470.

\bibitem{bruns1887}
Bruns, H. (1887).
\"{U}ber die Integrale des Vielk\"{o}rper-Problems.
\textit{Acta Mathematica}, \textbf{11}, 25--96.

\bibitem{burrau1913}
Burrau, C. (1913).
Numerische Berechnung eines Spezialfalles des Dreik\"{o}rperproblems.
\textit{Vierteljahrsschrift der Astronomischen Gesellschaft}, \textbf{48}, 135--136.

\bibitem{chenciner2000}
Chenciner, A., \& Montgomery, R. (2000).
A remarkable periodic solution of the three-body problem in the case of equal masses.
\textit{Annals of Mathematics}, \textbf{152}(3), 881--901.

\bibitem{chernin1994}
Chernin, A.\,D., Valtonen, M.\,J., \& Mikkola, S. (1994).
Chaos in a restricted four-body system: Kolmogorov--Sinai entropy and dynamics.
\textit{New Astronomy Reviews}, \textbf{38}, 251.

\bibitem{chernin2000}
Heinäm\"{a}ki, P., Lehto, H.\,J., Valtonen, M.\,J., \& Chernin, A.\,D. (2000).
Fractal distribution of points in the three-body problem: Kolmogorov--Sinai entropy.
\textit{Monthly Notices of the Royal Astronomical Society}, \textbf{310}, 811--822.

\bibitem{haller2001}
Haller, G. (2001).
Distinguished material surfaces and coherent structures in three-dimensional fluid flows.
\textit{Physica D: Nonlinear Phenomena}, \textbf{149}, 248--277.

\bibitem{heinamaki1998}
Heinäm\"{a}ki, P., Lehto, H.\,J., Valtonen, M.\,J., \& Chernin, A.\,D. (1998).
Three-body dynamics: intermittent chaos with strange attractor.
\textit{Monthly Notices of the Royal Astronomical Society}, \textbf{298}, 790--800.

\bibitem{hardy1979}
Hardy, G.\,H., \& Wright, E.\,M. (1979).
\textit{An Introduction to the Theory of Numbers} (5th ed.).
Oxford University Press.

\bibitem{heggie1975}
Heggie, D.\,C. (1975).
Binary evolution in stellar dynamics.
\textit{Monthly Notices of the Royal Astronomical Society}, \textbf{173}, 729--787.

\bibitem{heggie2003}
Heggie, D.\,C., \& Hut, P. (2003).
\textit{The Gravitational Million-Body Problem}.
Cambridge University Press.

\bibitem{hutbahcall1983}
Hut, P., \& Bahcall, J.\,N. (1983).
Binary--single star scattering. I -- Numerical experiments for equal masses.
\textit{Astrophysical Journal}, \textbf{268}, 319--341.

\bibitem{jacobi1842}
Jacobi, C.\,G.\,J. (1842).
Sur l'\'{e}limination des noeuds dans le probl\`{e}me des trois corps.
\textit{Astronomische Nachrichten}, \textbf{20}, 81--102.

\bibitem{kolmogorov1954}
Kolmogorov, A.\,N. (1954).
On conservation of conditionally periodic motions for a small change in Hamilton's function.
\textit{Doklady Akademii Nauk SSSR}, \textbf{98}, 527--530.

\bibitem{li2017}
Li, X., \& Liao, S. (2017).
More than six hundred new families of Newtonian periodic planar collisionless three-body orbits.
\textit{Science China Physics, Mechanics \& Astronomy}, \textbf{60}, 129511.
[arXiv:1705.00527]

\bibitem{laskar1990}
Laskar, J. (1990).
The chaotic motion of the solar system: a numerical estimate of the size of the chaotic zones.
\textit{Icarus}, \textbf{88}, 266--291.

\bibitem{laskar1993}
Laskar, J. (1993).
Frequency analysis for multi-dimensional systems: global dynamics and diffusion.
\textit{Physica D: Nonlinear Phenomena}, \textbf{67}, 257--281.

\bibitem{lehto2008}
Lehto, H.\,J., Kotiranta, S., Valtonen, M.\,J., Heinäm\"{a}ki, P., Mikkola, S., \& Chernin, A.\,D. (2008).
Mapping the three-body system: decay time and reversibility.
\textit{Monthly Notices of the Royal Astronomical Society}, \textbf{388}, 965--978.

\bibitem{leimkuhler2004}
Leimkuhler, B., \& Reich, S. (2004).
\textit{Simulating Hamiltonian Dynamics}.
Cambridge University Press.

\bibitem{mcdonald1985}
McDonald, S.\,W., Grebogi, C., Ott, E., \& Yorke, J.\,A. (1985).
Fractal basin boundaries.
\textit{Physica D: Nonlinear Phenomena}, \textbf{17}, 125--153.

\bibitem{moeckel2023}
Moeckel, R. (2023).
Natural dynamical reduction of the three-body problem.
\textit{Celestial Mechanics and Dynamical Astronomy}, \textbf{135}, 10.

\bibitem{montgomery2015}
Montgomery, R. (2015).
The three-body problem and the shape sphere.
\textit{American Mathematical Monthly}, \textbf{122}(4), 299--321.

\bibitem{moser1962}
Moser, J. (1962).
On invariant curves of area-preserving mappings of an annulus.
\textit{Nachrichten der Akademie der Wissenschaften G\"{o}ttingen}, \textbf{2}, 1--20.

\bibitem{nekhoroshev1977}
Nekhoroshev, N.\,N. (1977).
An exponential estimate of the time of stability of nearly integrable Hamiltonian systems.
\textit{Russian Mathematical Surveys}, \textbf{32}(6), 1--65.

\bibitem{newton1687}
Newton, I. (1687).
\textit{Philosophi\ae\ Naturalis Principia Mathematica}.
Royal Society, London.

\bibitem{poincare1892}
Poincar\'{e}, H. (1892--1899).
\textit{Les M\'{e}thodes Nouvelles de la M\'{e}canique C\'{e}leste}, Vols.\ I--III.
Gauthier-Villars, Paris.

\bibitem{portegies2007}
Portegies Zwart, S., Belleman, R.\,G., \& Geldof, P.\,M. (2007).
High-performance direct gravitational N-body simulations on graphics processing units.
\textit{New Astronomy}, \textbf{12}, 641--650.

\bibitem{rein2012}
Rein, H., \& Liu, S.-F. (2012).
REBOUND: an open-source multi-purpose N-body code for collisional dynamics.
\textit{Astronomy \& Astrophysics}, \textbf{537}, A128.

\bibitem{stone2019}
Stone, N.\,C., \& Leigh, N.\,W.\,C. (2019).
A statistical solution to the chaotic, non-hierarchical three-body problem.
\textit{Nature}, \textbf{576}, 406--410.

\bibitem{suvakov2013}
{\v{S}}uvakov, M., \& Dmitra{\v{s}}inovi{\'c}, V. (2013).
Three classes of Newtonian three-body planar periodic orbits.
\textit{Physical Review Letters}, \textbf{110}, 114301.

\bibitem{sundman1912}
Sundman, K.\,F. (1912).
M\'{e}moire sur le probl\`{e}me des trois corps.
\textit{Acta Mathematica}, \textbf{36}, 105--179.

\bibitem{szebehely1967book}
Szebehely, V. (1967).
\textit{Theory of Orbits: The Restricted Problem of Three Bodies}.
Academic Press, New York.

\bibitem{szebehely1967}
Szebehely, V., \& Peters, C.\,F. (1967).
Complete solution of a general problem of three bodies.
\textit{Astronomical Journal}, \textbf{72}, 876--883.

\bibitem{tanikawa2000}
Tanikawa, K., \& Mikkola, S. (2000).
One-dimensional three-body problem via symbolic dynamics.
\textit{Chaos: An Interdisciplinary Journal of Nonlinear Science}, \textbf{10}(3), 649--657.

\bibitem{tanikawa2015}
Tanikawa, K., \& Mikkola, S. (2015).
Symbol sequences and orbits of the free-fall three-body problem.
\textit{Publications of the Astronomical Society of Japan}, \textbf{67}(6), 115.

\bibitem{trani2024}
Trani, A.\,A., Leigh, N.\,W.\,C., Vesperini, E., \& Fujii, M.\,S. (2024).
Isles of regularity in a sea of chaos amid the gravitational three-body problem.
\textit{Astronomy \& Astrophysics}, \textbf{689}, A24.

\bibitem{valtonen2016}
Valtonen, M., Anosova, J., Kholshevnikov, K., Myll\"{a}ri, A., Orlov, V., \& Tanikawa, K. (2016).
\textit{The Three-body Problem from Pythagoras to Hawking}.
Springer International Publishing.

\bibitem{webgpu2024}
W3C GPU for the Web Working Group (2024).
WebGPU specification.
\url{https://www.w3.org/TR/webgpu/}

\bibitem{yoshida1990}
Yoshida, H. (1990).
Construction of higher order symplectic integrators.
\textit{Physics Letters A}, \textbf{150}(5--7), 262--268.

\end{thebibliography}

\end{document}
